\section{Changed terms and recovered credit}\label{sec:bank-restructuring}

A borrower can receive easier terms without recovering its ability to pay.
A bank can also recover its expectation of payment without recovering all
the credit quality present at origination. These are different changes to
the loan. The following historical Australian examples use the retained
AASB~9 compilation for synthetic 2024--2025 reporting periods. Probabilities,
legal findings and credit assessments are supplied assumptions; taxes and
fees are excluded.

Start from cash100, a gross loan110, allowance55 and deposits180. Net assets155
leave equity$-25$. At the 2024 close, the bank replaces a payment121 due
one year later with110 at the same date. Assume the change preserves the
existing asset's contractual rights and does not qualify for derecognition.
The original annual effective interest rate is10\%. Discounting the amended
contractual payment gives
\[
 G_{\mathrm{modified}}=\frac{110}{1.1}=100,
 \qquad \text{modification loss}=110-100=10.
\]
Debit modification expense10 and credit gross loan10. This calculation changes
the contractual carrying amount before considering expected defaults.
Paragraph5.4.3 requires that separation. Fees would require an additional
carrying-value adjustment and amortisation method
\citep[paragraph5.4.3]{aasb9dated}.

The bank separately expects full110 with probability$4/5$ and zero with
probability$1/5$. Expected cash88 has present value80, so the required
allowance is$100-80=20$. Debit the old allowance35 and credit impairment
reversal35. Cash100 and net loan80 now equal deposits180 and equity0.
Neither this reversal nor the new payment terms establish recovery from
credit impairment. For a loan that remains recognised, the comparison for a
significant increase in credit risk retains the original, unmodified
origination terms \citep[paragraph5.5.12 and B5.5.27]{aasb9dated}.

A distinct, ordered credit review at that same close changes the supplied
full-payment probability to$9/10$. Expected99 has present value90; allowance
falls20 to10, adding an impairment reversal10 and leaving equity10.
Compare two histories at this point. In one, reinstatement of the borrower's
operating licence and verified repayment capacity support a finding that the
asset is no longer credit-impaired. The improvement is objectively related to
an event after credit impairment. A significant increase in risk relative to
origination nevertheless remains. In the other history, the supplied credit
assessment still finds impairment despite the same improved cash expectation.
A forecast mean alone does not decide between these findings.

The distinction changes the next year's income presentation. For the cured
asset, paragraph5.4.2 returns subsequent interest to gross100: debit gross
loan10 and credit interest10. At maturity, the expected shortfall11 is no
longer discounted, so debit impairment expense1 and credit allowance1.
For the still-impaired asset, paragraph5.4.1(b) applies the original rate to
net90: debit gross loan10, credit interest9 and credit allowance1.
Both histories reach gross110, allowance11 and net99. Interest is10 in one
history and9 in the other, with the compensating difference in impairment
expense \citep[paragraphs5.4.1--2]{aasb9dated}.

The borrower then pays110, an outcome admitted by both supplied laws.
Debit cash110 and allowance11, credit gross110 and impairment reversal11.
Cash210 equals deposits180 plus equity30. Profit since the opening is55,
equal to110 received less the opening net loan55. Equal closing wealth
therefore cannot detect the wrong interest basis. Returning from lifetime
to twelve-month expected credit losses would require a separate reversal of
the significant-risk finding under paragraphs5.5.3--12.

Ending contractual rights requires another decision. Suppose a different
loan has gross100 and allowance20, and an enforceable agreement releases
all remaining rights for an actual payment85. There is no retained interest
or replacement instrument. Debit cash85 and allowance20, credit gross100
and gain5: consideration85 exceeds net carrying amount80 by5.
Without the release, the same payment cannot by itself establish full
derecognition \citep[paragraphs3.2.3 and3.2.12]{aasb9dated}.

For a transfer, the bank first identifies whether the unit is the whole
asset or an eligible part, then assesses transfer conditions, risks and
rewards, and control. A fixed-price repurchase can retain the risks rather
than establish a sale. The issuer-liability ten-percent comparison does not
answer these asset questions \citep[paragraphs3.2.2--7]{aasb9dated}.
A replacement asset needs fresh classification, initial measurement and a
purchased-or-originated credit-impairment assessment; a distressed replacement
is not automatically an ordinary new loan with twelve-month losses.
Partial transfers and continuing involvement require their own measurement
under paragraphs3.2.13--23 \citep[B5.5.25--26]{aasb9dated}.

\section{Collateral received after maturity}\label{sec:bank-late-collateral}

Collateral changes expected receipts through its amount, timing and
enforcement costs. Consider a separate loan advancing100 on1January2024
and requiring121 on1January2025. Its contractual half-year accumulation
factor is1.1, equivalent to21\% effective annual interest. This rate
convention differs from the preceding annual10\% loan.

The supplied initial law assigns probability$3/5$ to full payment at maturity
and$2/5$ to a default first observed within twelve months. In the default
branch, collateral receipts arrive on1July2025, one half-year after maturity.
Possible net receipts121,66.55 and12.10 have conditional probabilities
$1/4,1/2,1/4$, respectively. Their conditional mean is66.55. The bank has
supported reasonable prospects of fuller recovery before the sale;66.55 is
not a known upper limit. These probabilities are illustrative assumptions.

The expected loss compares the present value of contractual cash with that
of all expected receipts. The full-payment value is$121/1.1^2=100$; the
default-branch mean has value$66.55/1.1^3=50$. Thus
\[
 \operatorname{ECL}_0
 =100-\left(\frac35\,\frac{121}{1.1^2}
              +\frac25\,\frac{66.55}{1.1^3}\right)=20.
\]
The extra half-year belongs in the discounting. Omitting receipts after
maturity produces the wrong loss40. B5.5.55 expressly includes collateral
cash beyond maturity and deducts obtaining and selling costs; B5.5.43
limits twelve-month losses by the default date, not by the recovery date.
Here all possible defaults occur in the first year, so twelve-month and
lifetime losses coincide \citep[B5.5.43--44 and B5.5.55]{aasb9dated}.

Start with cash200, deposits180 and equity20. Advance100 and recognise
allowance20. On1July2024, first accrue gross interest10. Rolling the
original expectation forward gives net88 against gross110, so allowance
rises20 to22 through impairment expense2. A separately observed qualitative
credit-impairment event then establishes the collateral-only expectation:
$66.55/1.1^2=55$. Increase allowance22 to55 through impairment expense33.
The event is an admitted financial-distress finding, not an invented missed
instalment before any instalment is due.

During the next half-year, original-rate interest on net55 is5.5.
Debit gross loan11, credit interest5.5 and allowance5.5. At contractual
maturity gross121 less allowance60.5 leaves net60.5. There is no contractual
post-maturity interest, unamortised fee or premium. The initial100-to121
difference has been fully amortised. Over the final half-year the expected
receipt value rises60.5 to66.55, giving net interest6.05. Under these terms,
gross remains121 and allowance falls60.5 to54.45: debit allowance6.05 and
credit interest6.05. This follows from the AppendixA definitions and
B5.5.33's gross-minus-discounted-receipts measurement. It is a local derivation
for these terms, not a rule that every overdue contract stops accruing
contractual interest \citep[AppendixA, paragraph5.4.1(b), B5.5.33]{aasb9dated}.

The middle recovery occurs: an enforcement agent obtains79.86, deducts
costs13.31 already included in the forecast, and remits66.55.
Debit cash66.55 and credit gross loan66.55. Gross and allowance now both
equal54.45. At this sale, a separate supported exhaustion finding establishes
that no further recovery is reasonably expected. Paragraph5.4.4 requires
write-off: debit allowance54.45 and credit gross54.45. An earlier such
finding would require an earlier write-off; settlement is not a general
waiting rule \citep[paragraph5.4.4]{aasb9dated}.

Final cash166.55 equals deposits180 plus equity$-13.45$. Total interest21.55
less net impairment expense55 gives loss33.45, also equal to100 advanced
less66.55 received. The unpaid legal claim54.45 remains in the contract
history because no release has been supplied. Accounting write-off does not
by itself waive that claim. The bank never recognised the agent's collateral
as a separate owned asset, and never counted its proceeds a second time.
